\documentclass{article} % For LaTeX2e
\usepackage{iclr2026_conference,times}

% Optional math commands from https://github.com/goodfeli/dlbook_notation.
\input{math_commands.tex}

\usepackage{microtype}
\usepackage{graphicx}
\usepackage{amsmath,amssymb}
\usepackage{booktabs}
\usepackage{longtable}
\usepackage{hyperref}
\usepackage{url}
\usepackage{xspace}
\usepackage{listings}

\newcommand{\invarlock}{\textsc{InvarLock}\xspace}
\newcommand{\schemaverify}{\texttt{schema-verify}\xspace}

\lstset{
  basicstyle=\ttfamily\footnotesize,
  columns=fullflexible,
  breaklines=true,
  breakatwhitespace=true,
  keepspaces=true
}

\title{A Machine-Readable Schema for Weight-Edit\\Evaluation Reports with First-Class Verifier Traces}

% Authors must not appear in the submitted version.
\author{Anonymous Authors}

%\iclrfinalcopy % Uncomment for camera-ready version, but NOT for submission.
\begin{document}
\maketitle

\begin{abstract}
Evaluations of weight edits such as quantization and pruning increasingly mix
statistical guard frameworks (e.g., \invarlock), general evaluation harnesses,
and task-specific external verifiers (e.g., HumanEval/MBPP execution).
Today these tools emit incompatible artifacts, making it difficult to align guard
evidence with verifier outcomes, reproduce end-to-end verification loops, or
build automated promotion gates with auditable inputs.
We propose a lightweight JSON Schema for \emph{verifier-carrying evaluation artifacts}
that links three evidence elements in a single machine-checkable object under a
versioned wrapper (\texttt{schema\_version}):
\texttt{guard\_evidence} (tool-native guard measurements, thresholds, and verdicts, including
\invarlock evaluation reports and optional \texttt{invarlock verify} outputs),
\texttt{verifier\_traces} (external verifier results together with a trace contract covering prompt
hashes, harness identity and sandbox constraints, resource limits, and decoding parameters),
and \texttt{provenance} (model, tokenizer, edit configuration, and tool versioning).
A companion \schemaverify utility validates structural conformance and cross-block consistency,
enabling downstream automation to treat the artifact as a reliable unit of evidence in
verification-aware pipelines.
We implement serializers for \invarlock artifacts and code-execution verifiers and demonstrate
cross-tool queries on a paired dataset spanning multiple models and edit families, including cases
where guard evidence and verifier outcomes disagree.
We release the schema, serializers, validator, and a queryable demo corpus as open-source building
blocks for composing soft evidence with hard verification in model-edit workflows.
\end{abstract}

\section{Introduction}
When a language model is quantized, pruned, or otherwise weight-edited, evaluation evidence is
produced at three layers that rarely interoperate:
(i) \textbf{guard frameworks} (e.g., \invarlock) emit diagnostic measurements and gate flags,
(ii) \textbf{evaluation harnesses} define prompt sets and generation parameters, and
(iii) \textbf{external verifiers} (e.g., HumanEval/MBPP code-execution sandboxes) provide the
decisive outcome signal: did the edited model's outputs satisfy the verifier?
This three-layer structure aligns with VerifAI's emphasis on verifier-grounded outcomes:
verifier pass/fail is the headline metric, while guard evidence is auxiliary and diagnostic.
In practice, the layers emit incompatible artifacts with missing provenance; even basic
comparability questions become hard to answer automatically: are two pass@1 numbers computed over
the same prompt set, under the same harness revision and sandbox constraints, with the same
decoding parameters and seed?

Existing formats do not close this gap.
Model cards and leaderboards rarely encode verifier determinism contracts \citep{mitchell2019model_cards}.
Dataset metadata schemas such as Croissant improve dataset provenance but do not pin verifier-run
provenance \citep{akhtar2024croissant}.
Supply-chain frameworks like SLSA provide content-addressing primitives that inspire our design but
target software artifacts, not model-evaluation evidence \citep{slsa2023spec}.

\paragraph{Why not existing artifact and tracing standards?}
Experiment trackers (e.g., MLflow, W\&B) log metrics and artifacts but do not enforce a closed
schema with per-case verifier verdicts, prompt-set digesting rules, or code-execution sandbox
contracts.
SARIF \citep{sarif2020} standardizes static-analysis results but not generative-model evaluation.
OpenTelemetry \citep{opentelemetryspec} traces capture service-level spans, not the determinism
contract needed to compare two verifier runs.
The gap we fill is specific: a strict, machine-validated object that binds auxiliary guard evidence
to first-class verifier traces under a shared determinism contract, with per-case verdicts and
prompt-set integrity built in.

We address this with a \textbf{verifier-carrying artifact}: a strict JSON Schema that co-locates
hash-linked guard evidence with first-class verifier traces under shared provenance.
The verifier trace's contract pins comparability-critical conditions (prompt-set digests, harness
identity, model/tokenizer revisions, decoding parameters, and execution sandbox limits).
\schemaverify checks structural conformance and contract-driven cross-block consistency, without
claiming formal guarantees about model correctness or verifier soundness.

The resulting artifact supports verifier-centered workflows such as automated promotion gates and
escalation policies (escalate from cheap guard screens to expensive verifier runs when evidence
indicates risk).

\textbf{Contributions.}
(1) A verifier-trace schema and contract for auditable external verification runs.
(2) A wrapper schema that packages auxiliary guard evidence and first-class verifier traces in one
queryable object.
(3) A validator that enforces structure and cross-block consistency.

\section{Artifact Schema Overview}
The verifier-carrying artifact is a single JSON object
(\texttt{schema\_version} const \path{verifier_carrying_artifact.v1}) intended to be the unit of exchange for
weight-edit evaluation evidence.
It is intentionally strict (\texttt{additionalProperties=false}) on core objects so typos and
logging drift fail validation instead of silently passing.
It has three required evidence blocks under a single provenance envelope:
\texttt{provenance}, \texttt{guard\_evidence}, and \texttt{verifier\_traces[]} (minItems=1).

In v1, \texttt{guard\_evidence} includes \texttt{guard\_evidence.invarlock.evaluation\_report} as a
hash-linked file reference (with optional \texttt{embedded} for inline inclusion), plus optional
\texttt{guard\_evidence.invarlock.verify} output (\texttt{ok}, \texttt{profile}, \texttt{errors}, and
optional embedded \texttt{verify\_json} from \texttt{invarlock verify --json}).
The wrapper intentionally treats the \invarlock evaluation report as its own artifact (referenced
by digest), while the verifier trace is the primary verifier-first outcome record.
Future versions can add peer \texttt{guard\_evidence} blocks for other guard frameworks without
changing the verifier-trace contract.

\section{Verifier Trace Contract}
The trace contract is the minimal, machine-checkable record binding an external verifier outcome
(e.g., HumanEval/MBPP pass@1 \citep{chen2021humaneval,austin2021mbpp}) to the concrete
conditions under which it was produced.
If the contract fields match between two traces, the runs are \emph{intended to be comparable}; when
outcomes differ despite matching contracts, the artifact makes the mismatch diagnosable rather than
silent.

The v1 trace records:
(i) verifier identity (\texttt{verifier.name}, \texttt{verifier.kind}, and harness identity with at
least one of version, git commit, or container image),
(ii) prompt-set integrity (dataset identifiers plus ordered prompt ids and hashes and a dataset-level
\texttt{digest\_sha256} computed over canonical JSON bytes; \texttt{embedded} vs.\ \texttt{hash\_only}
modes),
(iii) determinism-oriented generation parameters (model/tokenizer id+revision, decoding parameters,
seed, and max token limits),
and (iv) execution sandbox bounds (required when \texttt{verifier.kind="code\_execution"}):
\texttt{timeout\_s}, \texttt{wall\_limit\_s}, \texttt{cpu\_limit}, \texttt{mem\_limit\_mb}, and schema-enforced
\texttt{network\_enabled=false}.

The contract is a trust boundary: it makes comparability and auditability machine-checkable, but it
does not prove determinism across platforms.
Hidden nondeterminism (GPU scheduling, floating-point ordering) can still cause outcome differences;
the contract ensures these differences are diagnosable because the controlled conditions are
recorded explicitly.

\section{Schema Verification (Validator)}
The \schemaverify utility validates:
(i) structural conformance against Draft 2020-12 JSON Schema \citep{jsonschema2020_12} for both the
wrapper artifact and each
verifier trace, and
(ii) cross-block consistency for contract-defined invariants, including prompt-set digest
recomputation, case-id alignment between prompt sets and results, and (optionally) on-disk SHA-256
verification when local paths are available.
These checks establish evidence integrity and joinability, not correctness.

\section{Demonstration Corpus and Example Queries}
We include a queryable demonstration summary,
\texttt{s1\_demo.summary.jsonl} (see Appendix~\ref{app:file-inventory}),
that illustrates how verifier-carrying artifacts enable joins and audit queries across tools.

\paragraph{Demo scope.}
The demo comprises 11 verifier-carrying artifacts, each paired with 2 verifier traces (HumanEval and
MBPP), yielding 22 joined rows.
It spans three base model families (Phi-2, DeepSeek-Coder-1.3B, and Qwen2.5-Coder-3B).
Two Phi-2 pruning variants appear as BYOE subject checkpoints; this distinction matters because the
trace contract pins the exact model id+revision used for generation.

Table~\ref{tab:demo} shows representative rows.

\begin{table}[t]
\centering
\scriptsize
\begin{tabular}{llllrrrr}
\toprule
Model & Edit & Verifier & $\Delta$pass@1 & PM & Spec & RMT & Spec.\ status \\
\midrule
DS-Coder-1.3B & INT8 RTN c=0.00 & HumanEval & +0.012 & \checkmark & \checkmark & \checkmark & capped \\
DS-Coder-1.3B & INT8 RTN c=0.00 & MBPP     & $-$0.006 & \checkmark & \checkmark & \checkmark & capped \\
Qwen-Coder-3B & INT8 RTN c=0.00 & HumanEval & $-$0.006 & \checkmark & \checkmark & \checkmark & capped \\
Qwen-Coder-3B & INT8 RTN c=0.00 & MBPP     & $-$0.002 & \checkmark & \checkmark & \checkmark & capped \\
Phi-2         & INT8 RTN c=0.00 & HumanEval & +0.098 & \checkmark & \checkmark & $\times$ & stable \\
Phi-2 (BYOE)  & prune 40\% & HumanEval & +0.220 & $\times$ & \checkmark & \checkmark & stable \\
\bottomrule
\end{tabular}
\caption{Representative demo rows (joined view), subset (6 of 22 rows). Full table in
Appendix~\ref{app:fulltable}. PM=primary-metric gate, Spec=spectral gate; c is clamp ratio.
``capped'' means spectral family caps were applied
(\texttt{caps\_applied>0}); it is not itself a failure. Gate outcomes are recorded separately.}
\label{tab:demo}
\end{table}

\paragraph{Queries.}
All example queries operate on the JSONL summary and demonstrate patterns that are hard to see in
per-tool logs.

\textbf{Q1: All main gates pass, but the verifier regresses (any amount).}
\begin{lstlisting}
jq -r '
  select(
    .gate_primary_metric_acceptable and
    .gate_invariants_pass and
    .gate_spectral_stable and
    .gate_rmt_stable and
    (.delta_pass_rate < 0)
  )
  | {artifact_path, invarlock_model_id, verifier_name,
     delta_pass_rate, verifier_n_total}
' s1_demo.summary.jsonl
\end{lstlisting}
\noindent\emph{On our demo, this surfaces mild MBPP regressions even when all gates pass
(e.g., $\Delta=-0.006$ on $N=500$). The contract field \texttt{verifier\_n\_total} makes the
significance question well-posed (a binomial 95\% CI is on the order of $\pm 0.03$).}

\textbf{Q2: Primary-metric gate fails, yet HumanEval improves materially.}
\begin{lstlisting}
jq -r '
  select(
    (not .gate_primary_metric_acceptable) and
    .verifier_name == "humaneval" and
    (.delta_pass_rate >= 0.10)
  )
  | {invarlock_model_id, invarlock_edit_name, delta_pass_rate,
     primary_metric_ratio_vs_baseline}
' s1_demo.summary.jsonl
\end{lstlisting}
\noindent\emph{This surfaces unnecessary rejections: edits the guard would block, yet the external
verifier improves (e.g., Phi-2 prune 40\% with $\Delta=+0.220$ on HumanEval).}

\textbf{Q3: Release-safe join via prompt digest (sorted before grouping).}
\begin{lstlisting}
jq -s '
  sort_by(.prompt_digest + ":" + .verifier_name)
  | group_by(.prompt_digest + ":" + .verifier_name)
  | map(select(length > 1))
  | .[] | map({invarlock_edit_name, delta_pass_rate, gate_spectral_stable})
' s1_demo.summary.jsonl
\end{lstlisting}
\noindent\emph{This demonstrates release-safe joinability: edits can be compared and grouped by
\texttt{prompt\_digest} without redistributing prompt text (via \texttt{hash\_only} prompt sets).}

\textbf{Q4: Caps applied but gate passes (capped $\neq$ failed).}
\begin{lstlisting}
jq -r '
  select(.spectral_status == "capped" and .gate_spectral_stable)
  | {invarlock_model_id, verifier_name, spectral_status,
     spectral_stability_score, gate_spectral_stable}
' s1_demo.summary.jsonl
\end{lstlisting}
\noindent\emph{This returns cases where \texttt{spectral\_status="capped"} but the gate still passes
(\texttt{gate\_spectral\_stable=true}), motivating distinct fields for status vs.\ verdict.}

\section{Limitations}
This work standardizes \emph{evidence}, not guarantees.
The validator checks structural well-formedness and contract-driven consistency; it cannot certify
verifier soundness or harness correctness without rerunning the underlying pipeline.
The trace contract records conditions needed for comparability but does not eliminate platform
dependence or hidden nondeterminism; the goal is auditability.
The demonstration corpus is small and meant to illustrate joins and queries; larger-scale analyses
of guard--verifier alignment are out of scope here.
The v1 \texttt{guard\_evidence} block is \invarlock-specific, while the \texttt{verifier\_trace} schema
is tool-agnostic; other guard frameworks can be supported via peer blocks.
Although the trace schema supports multiple verifier kinds (proof checking, SMT, static analysis),
our demo instantiates code-execution verifiers only.

\bibliographystyle{iclr2026_conference}
\bibliography{references}

\clearpage
\appendix
\section{Full Demonstration Table}\label{app:fulltable}
\small
% Auto-generated from research/verifai2_2026/examples/s1_demo.summary.jsonl
\setlength{\tabcolsep}{3pt}
\renewcommand{\arraystretch}{1.0}
\begin{longtable}{lllrcccc}
\toprule
Model & Edit & Ver. & $\Delta$ & PM & Spec & RMT & Spec.\ status\\
\midrule
\endfirsthead
\toprule
Model & Edit & Ver. & $\Delta$ & PM & Spec & RMT & Spec.\ status\\
\midrule
\endhead
DS-Coder-1.3B & INT8 RTN c=0.00 & HE & +0.012 & \checkmark & \checkmark & \checkmark & capped\\
DS-Coder-1.3B & INT8 RTN c=0.00 & MBPP & -0.006 & \checkmark & \checkmark & \checkmark & capped\\
DS-Coder-1.3B & INT8 RTN c=0.1 & HE & +0.293 & $\times$ & \checkmark & $\times$ & capped\\
DS-Coder-1.3B & INT8 RTN c=0.1 & MBPP & -0.126 & $\times$ & \checkmark & $\times$ & capped\\
DS-Coder-1.3B & INT8 RTN c=0.25 & HE & -0.238 & $\times$ & \checkmark & $\times$ & capped\\
DS-Coder-1.3B & INT8 RTN c=0.25 & MBPP & -0.126 & $\times$ & \checkmark & $\times$ & capped\\
Phi-2 & INT8 RTN c=0.00 & HE & +0.098 & \checkmark & \checkmark & $\times$ & stable\\
Phi-2 & INT8 RTN c=0.00 & MBPP & -0.004 & \checkmark & \checkmark & $\times$ & stable\\
Phi-2 & INT8 RTN c=0.25 & HE & -0.305 & $\times$ & \checkmark & $\times$ & stable\\
Phi-2 & INT8 RTN c=0.25 & MBPP & -0.032 & $\times$ & \checkmark & $\times$ & stable\\
Phi-2 & INT8 RTN c=0.5 & HE & -0.305 & $\times$ & \checkmark & $\times$ & capped\\
Phi-2 & INT8 RTN c=0.5 & MBPP & -0.032 & $\times$ & \checkmark & $\times$ & capped\\
Phi-2 (prune20) & Prune 20\% & HE & +0.037 & \checkmark & \checkmark & $\times$ & stable\\
Phi-2 (prune20) & Prune 20\% & MBPP & +0.000 & \checkmark & \checkmark & $\times$ & stable\\
Phi-2 (prune40) & Prune 40\% & HE & +0.220 & $\times$ & \checkmark & \checkmark & stable\\
Phi-2 (prune40) & Prune 40\% & MBPP & -0.002 & $\times$ & \checkmark & \checkmark & stable\\
Qwen-Coder-3B & INT8 RTN c=0.00 & HE & -0.006 & \checkmark & \checkmark & \checkmark & capped\\
Qwen-Coder-3B & INT8 RTN c=0.00 & MBPP & -0.002 & \checkmark & \checkmark & \checkmark & capped\\
Qwen-Coder-3B & INT8 RTN c=0.1 & HE & -0.238 & $\times$ & \checkmark & $\times$ & capped\\
Qwen-Coder-3B & INT8 RTN c=0.1 & MBPP & -0.404 & $\times$ & \checkmark & $\times$ & capped\\
Qwen-Coder-3B & INT8 RTN c=0.25 & HE & -0.238 & $\times$ & $\times$ & $\times$ & capped\\
Qwen-Coder-3B & INT8 RTN c=0.25 & MBPP & -0.404 & $\times$ & $\times$ & $\times$ & capped\\
\bottomrule
\end{longtable}


\clearpage
\section{Named File Inventory}\label{app:file-inventory}
To make file references concrete, we include a compact inventory of named files (by path) and
content digests (SHA-256 prefixes). Full digests are computed over raw file bytes.

\paragraph{Inventory (SHA-256 prefixes).}
\begin{center}
\scriptsize
% Auto-generated: file inventory (sha256 prefixes)
\setlength{\tabcolsep}{3pt}
\renewcommand{\arraystretch}{1.0}
\begin{tabular}{p{0.77\linewidth} l}
\toprule
File & SHA-256 prefix\\
\midrule
\path{research/verifai2_2026/specs/verifier_carrying_artifact.v1.schema.json} & \texttt{71a65f71c2f87f14}\\
\path{research/verifai2_2026/specs/verifier_trace.v1.schema.json} & \texttt{6c69d4eae828ee27}\\
\path{research/verifai2_2026/verifier_trace_contract.md} & \texttt{00ccddcc587aa582}\\
\path{scripts/verifai2_2026/schema_verify.py} & \texttt{caeb047b646042b9}\\
\path{scripts/verifai2_2026/verifier_trace_from_cases.py} & \texttt{a902f30993f7f3e4}\\
\path{scripts/verifai2_2026/cases_from_harness.py} & \texttt{566d6eec66dcb7a8}\\
\path{scripts/verifai2_2026/run_verifier_trace_pipeline.py} & \texttt{d8680da5bab256bf}\\
\path{research/verifai2_2026/examples/artifact.example.json} & \texttt{06ddda922bde77b0}\\
\path{research/verifai2_2026/examples/verifier_trace_humaneval.example.json} & \texttt{0e35a029364e7071}\\
\path{research/verifai2_2026/examples/verifier_trace_mbpp.example.json} & \texttt{1f5dec0923d4f288}\\
\path{research/verifai2_2026/examples/s1_demo.summary.jsonl} & \texttt{5e3024afa2e18474}\\
\bottomrule
\end{tabular}

\end{center}

\end{document}
